% !TEX root = ../main.tex
% Resultados. Nenhum numero digitado: todo valor vem de uma macro \F... de fatos.tex,
% gerada por artigo/tools/folha_de_fatos.py. O conferidor (folha_de_fatos.py --conferir)
% recusa digito solto neste arquivo.
% Duas camadas: a primeira frase de cada paragrafo diz em palavras o que se encontrou, e os
% numeros vem entre parenteses como prova. Cada conceito estatistico e traduzido uma vez, na
% primeira vez que aparece. Cada RQ fecha com uma caixa de resposta (\rqanswer, em main.tex).

\section{Results}
\label{sec:results}

Unless stated otherwise, results refer to the deterministic branch (\FarquivosDet{}
production files) and to density, the primary outcome. Table~\ref{tab:rq1} reports the
three outcomes in both branches.

\begin{table*}[t]
  \centering
  \caption{Spearman's $\rho$ between each outcome and the aggregate label. The partial
  $\rho$ removes the effect of the member count of the linked tests.}
  \label{tab:rq1}
  \begin{tabular}{lrrrrr}
    \toprule
    Outcome & deterministic $\rho$ [95\% CI] & deterministic $p$ & deterministic partial $\rho$
      & extended $\rho$ & extended $p$ \\
    \midrule
    density & \FrhoDensidadeDet{} [\FicDensidadeBaixo{}, \FicDensidadeAlto{}] & \FpDensidadeDet{}
      & \FparcialDensidade{} & \FrhoDensidadeAmp{} & \FpDensidadeAmp{} \\
    raw & \FrhoBrutoDet{} [\FicBrutoBaixo{}, \FicBrutoAlto{}] & \FpBrutoDet{}
      & \FparcialBruto{} & \FrhoBrutoAmp{} & \FpBrutoAmp{} \\
    distinct & \FrhoDistintosDet{} [\FicDistintosBaixo{}, \FicDistintosAlto{}] & \FpDistintosDet{}
      & \FparcialDistintos{} & \FrhoDistintosAmp{} & \FpDistintosAmp{} \\
    \bottomrule
  \end{tabular}
\end{table*}

\subsection{RQ1: Perceived Severity and Test Smell Density}
\label{sec:res-rq1}

Classes that reviewers judged more severe are not tested by tests with more test smells
per member. The correlation is close to zero: $\rho$ = \FrhoDensidadeDet{}, where zero
means no relation and one means that the files would be ranked in the same order by
severity and by density. The $p$-value, the probability of a correlation at least this far
from zero if there were no relation, is \FpDensidadeDet{}, above the usual threshold of
$0.05$.

This does not show that the relation is absent. The confidence interval, the range of
correlations compatible with the data, goes from \FicDensidadeBaixo{} to
\FicDensidadeAlto{} at 95\%. To ask whether the correlation is smaller than
\FmargemEquivalencia{}, the conventional threshold of a small correlation, we use an
equivalence test: at
90\% the interval ends at \FmaiorEfeitoDet{}, slightly above that threshold, so the test
cannot conclude that the correlation is smaller than it ($p$ = \FpTostDensidadeDet{}). The
data rule out correlations larger than \FmaiorEfeitoDet{}, but not a small positive one
below that value.

Counting test smells without dividing by test size gives a different picture: the raw
count does grow with severity ($\rho$ = \FrhoBrutoDet{}, $p$ = \FpBrutoDet{}). Part of this
comes from size. Larger tests have more smells ($\rho$ = \FrhoBrutoTamanho{} between the
raw count and the member count of the tests), and classes judged more severe have slightly
larger tests ($\rho$ = \FrhoRotuloMetodos{}). The partial correlation, which removes the
effect of test size, drops from \FrhoBrutoDet{} to \FparcialBruto{} for the raw count. For
density, which already divides by size, it stays at \FparcialDensidade{}. The number of
distinct smells lies between the two ($\rho$ = \FrhoDistintosDet{}, $p$ =
\FpDistintosDet{}) and, like the raw count, drops when test size is removed (partial
$\rho$ = \FparcialDistintos{}).

Comparing groups of files leads to the same picture. Files whose label the majority of
reviewers confirmed (\FnConfirmados{}) have more raw occurrences than clean files
(\FnLimpos{}), but not a higher density. We measure the difference between two groups with
the rank-biserial correlation, which reads like $\rho$: zero means the groups do not
differ. It is \FdeltaBrutoExtremos{} for the raw count ($p$ = \FpajBrutoExtremos{}) and
\FdeltaDensidadeExtremos{} for density ($p$ = \FpajDensidadeExtremos{}); both $p$-values
are adjusted by Benjamini--Hochberg for the three outcomes compared. Confirmed files and
files flagged but not confirmed (\FnDisputados{}) differ by \FdeltaDisputados{} in
density, but this comparison cannot answer the question: with groups of this size, only a
difference of \FdeltaDetectavelDisputados{} or more would be detected with 80\% power,
i.e., in eight out of ten samples.

The result holds under the two robustness checks. Files of the same repository may
resemble each other; resampling whole repositories instead of files, the 95\% interval of
density goes from \FbootBaixoDet{} to \FbootAltoDet{}, close to the original. Removing the
files affected by the nested class defect of the detector (\FnSemAninhadaDet{} files
remain), the correlation is \FrhoSemAninhadaDet{} ($p$ = \FpSemAninhadaDet{}).

\rqanswer{RQ1}{We found no relation between the severity that reviewers perceived in a
production class and the density of test smells in its tests. The data rule out a relation
of moderate or large size, but not a small one. Counting test smells without dividing by
test size produces a relation, part of which comes from test size.}

\subsection{RQ2: Individual Test Smells and Smell Pairs}
\label{sec:res-rq2}

No individual test smell, and no pair of code smell and test smell, is related to
perceived severity once we account for the number of tests. When many tests are run, some
fall below $p = 0.05$ by chance alone; the Benjamini--Hochberg correction adjusts each
$p$-value for that. Of the \FtestesPorSmell{} tests per test smell (each smell as raw
count and as density, on all files and again without the files with nested classes),
\FnominaisPorSmell{} fall below $0.05$ before the correction, and after it the smallest $p$-value is
\FmenorPajPorSmell{}. In the \FtestesGrade{} tests for pairs of code smell and test smell,
the smallest corrected $p$-value is \FmenorPajGrade{}. The uncorrected results are listed
in the replication package.

\rqanswer{RQ2}{No individual test smell explains the absence of a relation in RQ1. We
found no smells with opposite effects that cancel out in the aggregate, at least none that
our tests could detect.}

\subsection{RQ3: Normalization and Binding}
\label{sec:res-rq3}

Each of the two measurement decisions changes the answer to RQ1.

\paragraph{Normalization}
Whether test smells are counted or divided by test size decides whether a relation
appears. As a raw count, they are related to perceived severity ($p$ = \FpBrutoDet{} for
the correlation, \FpajBrutoExtremos{} for confirmed versus clean files); as density, they
are not ($p$ = \FpDensidadeDet{} and \FpajDensidadeExtremos{}). As shown in RQ1, the raw
count carries the size of the test, and test size is itself related to the label. A study
that counts test smells without dividing by size measures, in part, test size.

\paragraph{Binding}
Linking each class to every test that references it also changes the answer, because most
of these tests are not dedicated to the class. In the extended branch (\FarquivosAmp{}
files), the correlation is close to zero for all three outcomes (Table~\ref{tab:rq1}).
With more files, the estimate is narrower than in the deterministic branch: at 90\% the
interval ends at \FmaiorEfeitoAmp{}, below the threshold of a small correlation, and
resampling repositories gives an interval from \FbootBaixoAmp{} to \FbootAltoAmp{}. In the
audit, however, the class is the target of the test in only \FprecisaoHumanaAmp{} of the
pairs of this branch, against \FprecisaoHumanaDet{} in the deterministic branch; the test
executes code of the class in \FprecisaoHumanaAmpAmplo{} of them. Most tests linked by a
static reference use the class without testing it.

\rqanswer{RQ3}{Both decisions change the result. Without dividing by test size, a relation
appears that partly reflects test size. Linking classes to every test that references
them yields a precise null result, but about tests that, for the most part, do not test
the class.}
